% Guia do professor — Capítulo 19: Dispersão de Poluentes
\section{Capítulo 19 --- Dispersão de Poluentes}

\subsection*{Visão geral e ritmo}

O Cap.~18 aplicado a um problema com consequência social imediata.
Para físicos é um banquete de métodos: equação do calor, método das
imagens, balanço de caixa e movimento browniano — cada um resolvendo
um pedaço do problema regulatório. \textbf{Ritmo: 2 aulas de 2\,h
(cabe em 1 aula $+$ estudo dirigido se o semestre apertar).}

\begin{itemize}
  \item \textbf{Aula 1} — pluma gaussiana (dedução + imagens),
        máximo no solo e o $1/H^2$, as três plumas da estabilidade
        (Casos~1--2).
  \item \textbf{Aula 2} — modelo de caixa e episódios (Caso~3),
        fumigação, passeio aleatório e Taylor (Caso~4).
\end{itemize}

\subsection*{Objetivos de aprendizagem}

\begin{enumerate}
  \item deduzir a pluma gaussiana e aplicar a fonte-imagem;
  \item localizar e escalar o máximo de concentração no solo;
  \item associar forma de pluma a estabilidade e hora do dia;
  \item usar o modelo de caixa e o fator de ventilação para explicar
        episódios;
  \item conectar passeio aleatório, teorema de Taylor e difusividade
        turbulenta.
\end{enumerate}

\subsection*{Pré-requisitos a reativar}

CLA completa (Cap.~18: $z_i$, estabilidade, TKE, $T_L$); perfis e
inversões (Cap.~5); método das imagens (eletrostática!); equação do
calor; ventilação anticiclônica (Cap.~12).

\subsection*{Armadilhas conceituais frequentes}

\begin{itemize}
  \item \textbf{``Chaminé alta resolve''}: reduz o pico local
        ($1/H^2$) e EXPORTA o problema — gancho histórico da chuva
        ácida.
  \item \textbf{``Episódio $=$ mais emissão''}: os episódios de
        inverno são queda do fator de ventilação com emissão igual —
        o Caso~3 desfaz o mito.
  \item \textbf{Pluma estável ``comportada''}: o fanning é uma bomba
        armada esperando a fumigação da manhã.
  \item \textbf{Gaussiana como lei exata}: vale para terreno plano,
        estacionário, sem química; é o LEGO básico dos modelos, não
        o edifício.
  \item \textbf{Máximo noturno de poluição atribuído ao tráfego}: o
        rush é às 18\,h; o pico de C é mais tarde porque a caixa
        encolhe (N3).
\end{itemize}

\subsection*{Demonstração de baixo custo}

Incenso em sala: perto do teto (estável, estratificado) vira fanning;
sobre um aquecedor, looping — as três plumas por R\$\,5. Dados
CETESB do último inverno: sobrepor série de MP$_{10}$ com $U$ e
discutir o fator de ventilação.

\subsection*{Problemas sugeridos do livro (exercícios do Cap.~19)}

Aquecimento: $C$ no solo para fontes dadas; classes de Pasquill.
Aprofundamento: $x_{max}$, caixa com números reais. Síntese: ``por
que os piores episódios de SP são em julho, se se dirige o ano
todo?''. \emph{Confira a numeração na sua cópia (v1.02b).}

\subsection*{Uso do notebook \texttt{cap19\_poluentes.ipynb}}

\begin{center}
\begin{tabular}{lll}
\hline
Caso & Conteúdo & Momento sugerido \\
\hline
1 & pluma gaussiana: campo e pico no solo & Aula 1 \\
2 & looping/coning/fanning & Aula 1 \\
3 & caixa da cidade: 3 ventilações & Aula 2 \\
4 & Langevin: balístico $\to$ difusivo & Aula 2 \\
\hline
\end{tabular}
\end{center}

O Caso~4 fecha o arco ``física estatística na atmosfera'' começado no
Cap.~18 — Einstein 1905 com turbulência no lugar de moléculas.

\subsection*{Exercícios complementares e soluções}

Nas notas: T1--T5 e N1--N4. Soluções em \texttt{cap19\_solucoes.ipynb}
(\emph{não distribuir}): T3 com \texttt{sympy}; N2 (fumigação) e N3
(máximo noturno) são os de maior rendimento conceitual. Pares:
T3$\leftrightarrow$N1, T4$\leftrightarrow$N3, T5$\leftrightarrow$N4.

\subsection*{Conexões com o resto do curso}

CLA e turbulência (Cap.~18, essencial); estabilidade (Cap.~5);
inversões e anticiclones (Caps.~5, 12); brisas redistribuindo poluição
(Cap.~17); deposição e nuvens (Caps.~6--7); química do ozônio e clima
(Cap.~21); os mesmos esquemas de advecção nos modelos (Cap.~20).
